\section{Introduction}\label{sec:intro}

When private specialists charge firms for collaboration support, why should a local government offer basic innovation intermediation free of charge? The European Digital Innovation Hubs network illustrates how publicly supported access to expertise and experimentation can coexist with market-priced services \citep{EuropeanCommissionEDIH,EDIHFAQ2026}. The issue is not merely public versus private ownership: what real fixed cost is justified to preserve basic access? Counts of supported firms or successful ventures cannot settle it, because participation and the regional incidence of venture returns both matter.

The question becomes less straightforward when research partners have complementary technical skills. Greater complementarity makes potential collaborations more valuable, but a fee-setting specialist can respond by charging more and serving a smaller, more selected group of founders. Offering free basic access alongside a paid premium then restores participation excluded by the private-only fee, while also changing the use of premium services. Meanwhile, successful ventures may displace incumbent profits or generate benefits outside the financing jurisdiction. A technological improvement can therefore increase the benefits recovered by public access while also worsening the local market return from the additional collaborations. The balance between these forces, rather than either effect in isolation, determines the change in the value of free access.

The constituent mechanisms are well established. Private matching fees can screen heterogeneous participants and alter both participation and composition \citep{DamianoLi2007,DamianoLi2008,Hagiu2004b}. Free public and paid private services may coexist with endogenous private fees \citep{HuaChenZhang2016,PeitzSato2026,LiuReshidiRivadeneyra2026}. Innovation spillovers and product-market rivalry also separate local and aggregate welfare \citep{DAspremontJacquemin1988,CremerMarchandPestieau1997}. We do not claim these ingredients as new. Instead, we ask how technical complementarity changes the maximum fixed cost of adding free basic access when the private intermediary reoptimises its fee, and why the resulting policy-value response can have the opposite sign from local returns to commercialisation. The question concerns a specific public-access margin, not a general theory of screening or an unrestricted comparison of policy instruments.

We examine two jurisdictions, each with a welfare-maximising government, a profit-maximising private intermediary and heterogeneous founders. Provision decisions precede private fees, participation, cross-region research and differentiated-product Cournot competition. For tractability, founder ability is uniformly distributed and perfectly aligned with willingness to pay for a premium that provides direct service value, not improved research productivity. Governments choose only whether to supply zero-priced basic access at a real fixed cost. The central comparison fixes full partner-region participation while allowing the local private fee to adjust. These restrictions are material to the result.

On the interior private-only branch, greater complementarity raises the optimal fee and cutoff, reducing participation while increasing the average ability of those served. Free basic access restores participation, partly displacing premium fee receipts. These are familiar selection and mixed-provision mechanisms; their welfare interaction is the focus here.

Proposition~1 identifies an intermediate interval for the local founder's share of venture profits: large enough for complementarity to strengthen private fee-induced exclusion, yet small enough for net local commercialisation incidence to remain negative. Within this interval, under the proposition's conditions, greater complementarity raises the gross public-access cost ceiling despite adverse regional venture returns. This positive slope is proved for sufficiently weak positive matching on a feasible complementarity interval with a positive success baseline; an exact-rational certificate verifies its signs over a continuous two-parameter region at the original finite matching intensity. The result does not establish that public provision is always desirable or that the four-player entry equilibrium changes.

Local and national cost ceilings differ by the welfare increment in the partner jurisdiction. This familiar externality is evaluated here with privately selected founders and regionally assigned venture returns. Three separately certified numerical environments exhibit two cases of public underprovision relative to the restricted national benchmark and one case of agreement. The change in the four-party configuration requires a narrow fixed-cost interval and disappears under the reported Beta$(2,2)$ ability alternative. These are equilibrium witnesses, not a general classification or calibrated programme effects.

The contribution is the precise condition under which private fee-induced exclusion reverses the implication of negative local commercialisation incidence for the \emph{change} in public-access value. It is not a new screening theorem, a claim that venture entry is generally welfare improving, or a justification for constructing physical innovation hubs.

Section~\ref{sec:model} presents the model, Section~\ref{sec:equilibrium} private pricing and Section~\ref{sec:propositions} public-cost ceilings. Sections~\ref{sec:welfare} and \ref{sec:numerical} report equilibrium examples and robustness; Sections~\ref{sec:empirical}, \ref{sec:literature} and \ref{sec:institutions} address interpretation and prior work. Section~\ref{sec:conclusion} concludes, and proofs appear in Online Resource~1.
